% =====================================================================
%  4. Resultados e discussão — EM ANDAMENTO
%  PRISMA 2020: seleção (itens 16a e 16b) e caracterização (item 17, parcial).
%  PENDENTE: consolidar os quantitativos por etapa e o fluxograma completo.
%  PENDENTE: desenvolver a discussão dos desafios, benefícios e perspectivas.
% =====================================================================
\section{Resultados e discussão}\label{sec:resultados}

\subsection{Seleção dos estudos}\label{sec:avaliacao}

As decisões documentadas para esta revisão compreendem nove estudos
aprovados (Tabela~\ref{tab:aprovados}), três estudos excluídos
(Tabela~\ref{tab:reprovados}) e duas publicações com texto integral não
recuperado (Tabela~\ref{tab:nao-recuperados}). O uso complementar do resumo
de Gomes et al.~\citeano{gomes2025}, descrito na
Seção~\ref{sec:fontes-complementares}, mantém a condição de texto integral
não recuperado desse estudo e não altera o total de estudos aprovados.

Os quantitativos das etapas de identificação e triagem e o fluxograma
completo da seleção ainda dependem da consolidação dos registros de busca.

{\footnotesize
\begin{longtable}{@{}>{\raggedright\arraybackslash}p{4.2cm}
                    p{\dimexpr\linewidth-4.2cm-2\tabcolsep\relax}@{}}
\caption{Estudos aprovados}\label{tab:aprovados}\\
\toprule
\textbf{Estudo} & \textbf{Justificativa} \\
\midrule
\endfirsthead
\caption[]{Estudos aprovados (continuação)}\\
\toprule
\textbf{Estudo} & \textbf{Justificativa} \\
\midrule
\endhead
\bottomrule
\endfoot
\cite{giamattei2024}\newline \emph{Monitoring tools for DevOps and microservices}
  & Analisa ferramentas de monitoramento em ambientes DevOps e microsserviços,
    contribuindo para a compreensão das práticas de observabilidade, métricas,
    logs e traces. \\
\addlinespace
\cite{plant2025}\newline \emph{Improving DevOps team performance through
  context-capability coalignment}
  & Investiga capacidades e desempenho de equipes DevOps em organizações
    públicas, estabelecendo conexão direta entre cultura DevOps e
    Administração Pública. \\
\addlinespace
\cite{azad2023}\newline \emph{DevOps critical success factors}
  & Identifica fatores críticos de sucesso na adoção de DevOps, fundamentando
    aspectos culturais, técnicos e organizacionais relevantes para o
    estudo. \\
\addlinespace
\cite{amaro2025}\newline \emph{Mapping DevOps capabilities to the software
  life cycle}
  & Relaciona capacidades DevOps às etapas do ciclo de vida do software,
    contribuindo para a fundamentação das práticas de automação,
    monitoramento e melhoria contínua. \\
\addlinespace
\cite{figueredo2026}\newline \emph{Challenges of Digital Services in Public
  Administration in the Era of Digital Transformation}
  & Identifica barreiras institucionais, organizacionais e tecnológicas à
    transformação digital na Administração Pública, contextualizando os
    desafios de adoção de novas práticas de engenharia de software. \\
\addlinespace
\cite{amaro2026}\newline \emph{Strategic initiatives in DevOps adoption}
  & Investiga a relação entre capacidades técnicas, métricas, processos e
    iniciativas estratégicas na adoção de DevOps, contribuindo para a
    discussão sobre integração entre cultura organizacional e práticas
    tecnológicas. \\
\addlinespace
\cite{linjama2026}\newline \emph{Continuous Deployment Adoption: Insights From
  a Public Sector Implementation}
  & Apresenta um estudo de caso sobre implantação de entrega contínua no
    setor público, oferecendo evidências sobre benefícios, desafios e
    particularidades da adoção de práticas DevOps na Administração
    Pública. \\
\addlinespace
\cite{kosinska2023}\newline \emph{Toward the Observability of Cloud-Native
  Applications}
  & Apresenta um mapeamento sistemático sobre observabilidade em aplicações
    cloud-native, contribuindo para a fundamentação conceitual, abordagens
    de implementação e perspectivas de observabilidade em sistemas
    distribuídos. \\
\addlinespace
\cite{alongi2022}\newline \emph{Event-sourced, observable software
  architectures}
  & Aborda a observabilidade como propriedade arquitetural de software,
    apresentando princípios e experiências práticas que complementam a
    fundamentação técnica sobre sistemas observáveis. \\
\end{longtable}
}

\begin{table}[H]
\centering
\caption{Estudos excluídos}\label{tab:reprovados}
\footnotesize
\begin{tabularx}{\linewidth}{@{}>{\raggedright\arraybackslash}p{4.2cm}
                                >{\raggedright\arraybackslash}p{2.2cm}X@{}}
\toprule
\textbf{Estudo} & \textbf{Decisão} & \textbf{Justificativa} \\
\midrule
\cite{gall2022}\newline \emph{Taking DevOps mainstream}
  & Excluído
  & Duplicidade de tratamento direto dos fundamentos conceituais de DevOps já
    escolhidos em outro artigo. \\
\addlinespace
\cite{hrusto2021}\newline \emph{Closing the feedback loop in DevOps through
  autonomous monitors in operations}
  & Excluído
  & Investiga uma solução específica de monitoramento autônomo e filtragem de
    alertas, em vez de uma análise abrangente da adoção de
    observabilidade. \\
\addlinespace
\cite{borges2024}\newline \emph{Informed and assessable observability design
  decisions in cloud-native microservice applications}
  & Excluído
  & Concentra-se em experimentos e decisões técnicas de observabilidade com
    injeção de falhas, um recorte arquitetural especializado. \\
\bottomrule
\end{tabularx}
\end{table}

\begin{table}[H]
\centering
\caption{Publicações com texto integral não recuperado}\label{tab:nao-recuperados}
\footnotesize
\begin{tabularx}{\linewidth}{@{}>{\raggedright\arraybackslash}p{4.2cm}X@{}}
\toprule
\textbf{Estudo} & \textbf{Situação e uso no artigo} \\
\midrule
\cite{gomes2025}\newline \emph{A systematic mapping study on observability of
  microservices-based applications}
  & Texto integral não recuperado. O resumo editorial é utilizado apenas
    como referência conceitual complementar
    (Seção~\ref{sec:fontes-complementares}); o estudo não integra o conjunto
    dos nove estudos aprovados. \\
\addlinespace
\cite{kumar2025}\newline \emph{A systematic literature review for
  investigating DevOps metrics to implement in software development
  organizations}
  & Texto integral não recuperado; o estudo não integra o conjunto dos
    estudos aprovados. \\
\bottomrule
\end{tabularx}
\end{table}

\FloatBarrier
